\section{Billed- og materialebilag: Træningsmester}
\label{app:traeningsmester-billedbilag}

\sloppy

Dette bilag samler udvalgte billeder, figurer og tabeller om Træningsmester pr. 28. april 2026. Bilaget omfatter 13 aktuelle iOS 26.4-screenshots, 7 arkitektur- og dataflowfigurer med supplerende ER-udsnit og et kondenseret udvalg af use cases, krav, SQL-/RLS-overblik, procesmodeller og anonymiseret beta-evidens. Ældre screenshots fra tidligere UI-versioner er ikke medtaget, fordi de ikke bør bruges som dokumentation for den aktuelle appversion.

Figurerne og tabellerne er udvalgt, så argumentationen kun henviser til materiale, der er indsat i bilaget. De supplerende tabeller er kondenserede og anonymiserede; private rådata, rå kommunikation og personhenførbare detaljer er ikke kopieret ind i bilaget.

\newcolumntype{Y}{>{\raggedright\arraybackslash}X}

\makeatletter
\setlength{\@fptop}{0pt}
\setlength{\@fpsep}{12pt plus 2pt minus 2pt}
\setlength{\@fpbot}{0pt plus 1fil}
\makeatother

\newcommand{\tmscreenshot}[3]{%
\begin{figure}[p]
    \centering
    \includegraphics[width=0.95\textwidth,height=0.74\textheight,keepaspectratio]{#1}
    \caption{#2}
    \label{#3}
\end{figure}
}

\newcommand{\tmdiagram}[3]{%
\begin{figure}[p]
    \centering
    \includegraphics[width=\textwidth,height=0.82\textheight,keepaspectratio]{#1}
    \caption{#2}
    \label{#3}
\end{figure}
}

\newcommand{\tmlandscapediagramwithsection}[5]{%
\clearpage
\begin{landscape}
\subsection{#1}
\vspace*{\fill}
\begin{center}
    \includegraphics[width=0.96\linewidth,height=0.62\textheight,keepaspectratio,trim=#3,clip]{#2}
    \captionof{figure}{#4}
    \label{#5}
\end{center}
\vspace*{\fill}
\end{landscape}
\clearpage
}

\newcommand{\tmlandscapediagram}[4]{%
\begin{landscape}
\vspace*{\fill}
\begin{center}
    \makebox[\linewidth][c]{\includegraphics[width=1.04\linewidth,height=0.72\textheight,keepaspectratio,trim=#2,clip]{#1}}
    \captionof{figure}{#3}
    \label{#4}
\end{center}
\vspace*{\fill}
\end{landscape}
\clearpage
}

\newcommand{\tmlandscapewidediagram}[4]{%
\begin{landscape}
\vspace*{\fill}
\begin{center}
    \makebox[\linewidth][c]{\includegraphics[width=1.08\linewidth,height=0.76\textheight,keepaspectratio,trim=#2,clip]{#1}}
    \captionof{figure}{#3}
    \label{#4}
\end{center}
\vspace*{\fill}
\end{landscape}
\clearpage
}

\newcommand{\tmlandscapeerdetail}[4]{%
\begin{landscape}
\vspace*{\fill}
\begin{center}
    \makebox[\linewidth][c]{\includegraphics[width=1.10\linewidth,height=0.79\textheight,keepaspectratio,trim=#2,clip]{#1}}
    \captionof{figure}{#3}
    \label{#4}
\end{center}
\vspace*{\fill}
\end{landscape}
\clearpage
}

\newcommand{\tmtrackercompletiondiagram}{%
\begin{landscape}
\vspace*{\fill}
\begin{center}
\resizebox{0.96\linewidth}{!}{%
\begin{tikzpicture}[
    >={Latex[length=3mm]},
    node distance=12mm,
    tmuser/.style={draw=blue!70!black, fill=blue!8, rounded corners=2pt, align=center, minimum width=30mm, minimum height=13mm, font=\footnotesize},
    tmdata/.style={draw=orange!80!black, fill=orange!10, rounded corners=2pt, align=center, minimum width=34mm, minimum height=13mm, font=\footnotesize},
    tmdecision/.style={diamond, aspect=2.1, draw=green!45!black, fill=green!8, align=center, inner sep=1.5pt, minimum width=24mm, minimum height=15mm, font=\footnotesize},
    tmarrow/.style={->, line width=0.8pt}
]
\node[tmuser] (start) at (0,0) {Bruger starter\\dagens workout};
\node[tmdecision] (tracker) at (3.7,0) {Tracker\\slået til?};
\node[tmdata] (logging) at (7.7,1.25) {WorkoutTracker\\logger sæt\\i trackerlog};
\node[tmdata] (notracker) at (7.7,-1.25) {No-tracker workout\\viser dagens\\øvelser};
\node[tmuser] (finish) at (11.8,0) {Bruger afslutter\\workout};
\node[tmdata] (completion) at (15.9,0) {WorkoutCompletion\\Service skriver\\completion-event};
\node[tmdecision] (cycle) at (20.2,0) {Aktiv training\\cycle runtime?};
\node[tmdata] (rpc) at (24.6,1.25) {RPC vælger\\næste slot};
\node[tmdata] (planflow) at (24.6,-1.25) {Planflow bevarer\\eller flytter\\normal næste dag};
\node[tmuser] (next) at (28.9,0) {Home, Watch\\og Tracker viser\\næste handling};

\draw[tmarrow] (start) -- (tracker);
\draw[tmarrow] (tracker) -- node[above, font=\scriptsize] {Ja} (logging);
\draw[tmarrow] (tracker) -- node[below, font=\scriptsize] {Nej} (notracker);
\draw[tmarrow] (logging) -- (finish);
\draw[tmarrow] (notracker) -- (finish);
\draw[tmarrow] (finish) -- (completion);
\draw[tmarrow] (completion) -- (cycle);
\draw[tmarrow] (cycle) -- node[above, font=\scriptsize] {Ja} (rpc);
\draw[tmarrow] (cycle) -- node[below, font=\scriptsize] {Nej} (planflow);
\draw[tmarrow] (rpc) -- (next);
\draw[tmarrow] (planflow) -- (next);
\end{tikzpicture}%
}
\captionof{figure}{Opsummering af tracker-on/off, completion-event og progression i træningscyklus.}
\label{fig:tm-tracker-completion-summary}
\end{center}
\vspace*{\fill}
\end{landscape}
\clearpage
}

\subsection{App-screenshots}

\tmscreenshot{Materiale/traeningsmester-2026-04-28/02_figurer_og_screenshots/app_screenshots/aktuel_ios26/ios26-01-home-stale-session-testkonto-2026-04-28.jpg}{Home-skærm med advarsel om en forældet tracker-session og mulighed for at fortsætte eller rydde sessionen.}{fig:tm-ios-home-stale-session}

\tmscreenshot{Materiale/traeningsmester-2026-04-28/02_figurer_og_screenshots/app_screenshots/aktuel_ios26/ios26-02-home-ready-testkonto-2026-04-28.jpg}{Home-skærm med aktuel træningshandling og næste planlagte træningspas.}{fig:tm-ios-home-ready}

\tmscreenshot{Materiale/traeningsmester-2026-04-28/02_figurer_og_screenshots/app_screenshots/aktuel_ios26/ios26-03-programmer-testkonto-2026-04-28.jpg}{Programoversigt med aktiv programstruktur og brugerens valgte progression.}{fig:tm-ios-programmer}

\tmscreenshot{Materiale/traeningsmester-2026-04-28/02_figurer_og_screenshots/app_screenshots/aktuel_ios26/ios26-04-oevelser-testkonto-2026-04-28.jpg}{Øvelseskatalog med søgning, filtrering og øvelsesliste i den aktuelle iOS-version.}{fig:tm-ios-oevelser}

\tmscreenshot{Materiale/traeningsmester-2026-04-28/02_figurer_og_screenshots/app_screenshots/aktuel_ios26/ios26-05-match-card-testkonto-2026-04-28.jpg}{Match-flow med swipe-kort til inspiration og valg af øvelser.}{fig:tm-ios-match-card}

\tmscreenshot{Materiale/traeningsmester-2026-04-28/02_figurer_og_screenshots/app_screenshots/aktuel_ios26/ios26-06-indstillinger-testkonto-2026-04-28.jpg}{Indstillingsoversigt med kontosektion og brugerens aktuelle app-tilstande.}{fig:tm-ios-indstillinger}

\tmscreenshot{Materiale/traeningsmester-2026-04-28/02_figurer_og_screenshots/app_screenshots/aktuel_ios26/ios26-07-konto-loginmetoder-testkonto-2026-04-28.jpg}{Kontoens loginmetoder i indstillingsflowet.}{fig:tm-ios-konto-loginmetoder}

\tmscreenshot{Materiale/traeningsmester-2026-04-28/02_figurer_og_screenshots/app_screenshots/aktuel_ios26/ios26-08-premium-abonnement-testkonto-2026-04-28.jpg}{Premium- og abonnementsvisning med appens betalingsrelaterede funktionsadgang.}{fig:tm-ios-premium-abonnement}

\tmscreenshot{Materiale/traeningsmester-2026-04-28/02_figurer_og_screenshots/app_screenshots/aktuel_ios26/ios26-09-historik-testkonto-2026-04-28.jpg}{Historikvisning med tidligere træningslog og gennemførte pas.}{fig:tm-ios-historik}

\tmscreenshot{Materiale/traeningsmester-2026-04-28/02_figurer_og_screenshots/app_screenshots/aktuel_ios26/ios26-10-tracker-active-testkonto-2026-04-28.jpg}{Aktiv tracker-session med planlagte sæt og løbende registrering under træning.}{fig:tm-ios-tracker-active}

\tmscreenshot{Materiale/traeningsmester-2026-04-28/02_figurer_og_screenshots/app_screenshots/aktuel_ios26/ios26-11-home-completion-delay-testkonto-2026-04-28.jpg}{Home-skærm efter gennemført pas med forsinket completion-feedback og fortrydelsesvindue.}{fig:tm-ios-home-completion-delay}

\tmscreenshot{Materiale/traeningsmester-2026-04-28/02_figurer_og_screenshots/app_screenshots/aktuel_ios26/ios26-12-home-next-training-testkonto-2026-04-28.jpg}{Home-skærm efter completion, hvor næste træningshandling er rykket frem.}{fig:tm-ios-home-next-training}

\tmscreenshot{Materiale/traeningsmester-2026-04-28/02_figurer_og_screenshots/app_screenshots/aktuel_ios26/ios26-13-home-current-xcode-mcp-qa-2026-04-28.jpg}{Home-skærm efter gennemløb med Dagens træning synlig.}{fig:tm-ios-home-current-qa}

\tmlandscapediagramwithsection{Arkitektur- og dataflowdiagrammer}{assets/bilag/diagrammer/system_context.pdf}{0bp 532bp 0bp 0bp}{Systemkontekst for Træningsmester på iOS, watchOS, Supabase og Apple-services.}{fig:tm-system-context}

\tmlandscapewidediagram{assets/bilag/diagrammer/data_model_er.pdf}{0bp 610bp 0bp 0bp}{ER-diagram for programmer, workouts, tracker, completion, settings og cyklusdata.}{fig:tm-data-model-er}

\tmlandscapeerdetail{assets/bilag/diagrammer/data_model_er.pdf}{0bp 610bp 360bp 0bp}{ER-diagram, venstre udsnit: bruger, programmer, workouts og deling.}{fig:tm-data-model-er-left}

\tmlandscapeerdetail{assets/bilag/diagrammer/data_model_er.pdf}{180bp 610bp 180bp 0bp}{ER-diagram, midterudsnit: settings, workouts, øvelser, trackerlog og completion-events.}{fig:tm-data-model-er-middle}

\tmlandscapeerdetail{assets/bilag/diagrammer/data_model_er.pdf}{360bp 610bp 0bp 0bp}{ER-diagram, højre udsnit: træner-/klientrelationer og træningscyklusdata.}{fig:tm-data-model-er-right}

\tmtrackercompletiondiagram

\tmlandscapediagram{assets/bilag/diagrammer/security_boundary_summary.pdf}{0bp 590bp 0bp 0bp}{Sikkerhedsgrænse mellem mobil klient, RLS og Edge Functions.}{fig:tm-security-boundary-summary}

\tmlandscapediagram{assets/bilag/diagrammer/watch_sync_sequence.pdf}{10bp 610bp 10bp 0bp}{Sekvens for iOS/watchOS auth-sync, token-cache og RLS-baseret workoutdata.}{fig:tm-watch-sync-sequence}

\tmlandscapediagram{assets/bilag/diagrammer/live_activity_idempotency_sequence.pdf}{15bp 515bp 15bp 5bp}{Sekvens for Live Activity intent-write, App Group snapshot og idempotent trackerlog.}{fig:tm-live-activity-idempotency-sequence}

\tmlandscapediagram{assets/bilag/diagrammer/import_ai_edge_boundary.pdf}{0bp 585bp 0bp 0bp}{Import/AI-grænse mellem klient, Edge Functions, OpenAI og Supabase.}{fig:tm-import-ai-edge-boundary}

\subsection{Materialegennemgang}

\begin{table}[htbp]
\small
\renewcommand{\arraystretch}{1.18}
\begin{tabularx}{\textwidth}{>{\raggedright\arraybackslash}p{0.23\textwidth} >{\raggedright\arraybackslash}p{0.28\textwidth} Y}
\textbf{Bilagsspor} & \textbf{Faglig funktion} & \textbf{Begrundelse} \\
\hline
Aktuelle app-screenshots & Dokumenterer den aktuelle brugerflade & Figurerne \ref{fig:tm-ios-home-stale-session}--\ref{fig:tm-ios-home-current-qa} viser centrale flows i iOS 26.4-versionen af Træningsmester. \\
Arkitektur- og dataflowdiagrammer & Dokumenterer systemets tekniske hovedstruktur & Figurerne \ref{fig:tm-system-context}--\ref{fig:tm-import-ai-edge-boundary} dækker systemkontekst, datamodel med forstørrede ER-udsnit, tracker/completion, RLS, watchOS, Live Activity og import/AI-grænser. \\
Containerarkitektur & Samler arkitekturen i komponentlag & Tabel \ref{tab:tm-container-architecture} gør hovedlagene læsbare som tekstligt overblik. \\
Trackerproces og autorisationsregler & Forklarer træningsflow og sikkerhedsmodel & Tabel \ref{tab:tm-tracker-flow-bpmn} og \ref{tab:tm-authorization-dcr} viser de centrale proces- og adgangsregler. \\
Use cases og kravsporbarhed & Kobler funktionalitet til implementering og test & Tabel \ref{tab:tm-use-cases} og \ref{tab:tm-krav-traceability} viser sammenhængen mellem brugerhandlinger, krav og implementeringsspor. \\
SQL, RLS og sikkerhed & Dokumenterer data- og adgangsmodellen & Tabel \ref{tab:tm-sql-domain-overblik}, \ref{tab:tm-sql-fields} og \ref{tab:tm-rls-matrix} beskriver datadomæner, idempotency-felter og sikkerhedsgrænser. \\
Beta- og TestFlight-spor & Dokumenterer iteration og brugerindsigt & Tabel \ref{tab:tm-testflight-timeline} og \ref{tab:tm-beta-feedback} samler anonymiserede observationspunkter uden direkte citater eller persondata. \\
Afgrænset råmateriale & Understøtter anonymisering og versionsklarhed & Private rådata, lukkede kommunikationsspor og ældre UI-optagelser er fravalgt af hensyn til samtykke, dataminimering og klarhed om den dokumenterede appversion. \\
\end{tabularx}
\caption{Gennemgang af bilagsmateriale og afgrænsning.}
\label{tab:tm-bilag-materialegennemgang}
\end{table}

\clearpage

\subsection{Use cases og krav}

\begin{table}[htbp]
\footnotesize
\renewcommand{\arraystretch}{1.12}
\begin{tabularx}{\textwidth}{>{\raggedright\arraybackslash}p{0.17\textwidth} >{\raggedright\arraybackslash}p{0.30\textwidth} Y}
\textbf{Use case} & \textbf{Kort formål} & \textbf{Faglig relevans} \\
\hline
UC-01 Opret konto og log ind & Supabase Auth, session restore og app-shell efter autentifikation. & Viser auth som forudsætning for brugerbundne træningsdata og RLS. \\
UC-02 Gennemfør onboarding & Valg af startvej, præferencer og eventuel programimport. & Understøtter analyse af første brugerrejse og friktion. \\
UC-03 Opret/rediger program & Oprette, hente, opdatere, aktivere og dele træningsprogrammer. & Dokumenterer programstruktur som kernefunktion. \\
UC-04 Tilføj træningsdag og øvelser & Knytte workouts og øvelser til programmet med sæt, reps, vægt og pause. & Knytter planlægning til den konkrete træningsfaglige struktur. \\
UC-05 Find og like øvelse & Søge, swipe og gemme liked/parrede øvelser. & Kobler øvelseskatalog, søgning og match-flow til brugerfeedback. \\
UC-06 Start workout med tracker & Logge sæt, styre pauser, opdatere Live Activity og gemme completion. & Central use case for fleksibel progression og lav-friktionslogging. \\
UC-07 Brug watchOS companion & Vise dagens træning og logge/afslutte workout fra Apple Watch. & Dokumenterer multiplatform-sporet uden runtime-screenshots. \\
UC-08 Aktiv træningscyklus & Prioritere cyklus-runtime over lineær planindekslogik. & Underbygger progression, deload og afvigelseshåndtering. \\
UC-09 Træner administrerer klient & Træner ser klienter, aktivitet og programlinks via servervalideret relation. & Afgrænser trænerfunktioner som teknisk udvidelse, ikke hovedproblem. \\
UC-10 Admin håndterer bruger eller aktivitet & Admin-center gennem Edge Function og server-side admin guard. & Underbygger governance, katalogkvalitet og sikkerhedsdesign. \\
UC-11 Importer program eller historik & Importjob, normalisering, preview og gem af program/historikkladder. & Viser AI/import som afgrænset server-side flow. \\
\end{tabularx}
\caption{Kondenseret use case-overblik for Træningsmester.}
\label{tab:tm-use-cases}
\end{table}

\begin{table}[htbp]
\footnotesize
\renewcommand{\arraystretch}{1.12}
\begin{tabularx}{\textwidth}{>{\raggedright\arraybackslash}p{0.23\textwidth} >{\raggedright\arraybackslash}p{0.19\textwidth} Y}
\textbf{Use case} & \textbf{Dækkede krav} & \textbf{Primære implementeringsspor} \\
\hline
UC-01 Opret konto og log ind & F-01, F-02, NF-01, NF-02 & Login, AppState, AppEnvironment, Supabase repositories. \\
UC-02 Gennemfør onboarding & F-02, F-11, F-12 & Onboardingmodeller, onboarding-flow og ProgramReviewService. \\
UC-03 Opret/rediger træningsprogram & F-03, F-04, NF-02, NF-03 & TrainingProgramView, AddProgramView og PlanRepository. \\
UC-04 Tilføj træningsdag og øvelser & F-04, F-05, F-06 & WorkoutRepository, TrackerRepository og workout exercise payloads. \\
UC-05 Find og like øvelse & F-06, F-07 & ExercisesView, MatchView og ExerciseSearchCatalog. \\
UC-06 Start workout med tracker & F-08, F-09, F-10, F-15 & WorkoutTrackerView, WorkoutCompletionService og Live Activity-spor. \\
UC-07 Brug watchOS companion & F-14, F-08, F-09, NF-07 & WatchWorkoutStore og WatchAuthSyncBridge. \\
UC-08 Aktiv træningscyklus & F-13, F-08, F-14 & TrainingCycleViews, TrainingCycleRepository og cycle RPC. \\
UC-09 Træner administrerer klient & F-16, NF-02, NF-08 & ClientsView, TrainerClientService og trainer-client Edge Function. \\
UC-10 Admin håndterer bruger & F-17, NF-01, NF-02 & SettingsView, admin-control-center og admin SQL functions. \\
UC-11 Importer program/historik & F-12, NF-08 & PremiumView, ProgramReviewService og import Edge Functions. \\
\end{tabularx}
\caption{Sporbarhed mellem use cases, krav og implementeringsspor.}
\label{tab:tm-krav-traceability}
\end{table}

\clearpage

\subsection{Arkitektur-, proces- og autorisationsmodeller}

\begin{table}[htbp]
\small
\renewcommand{\arraystretch}{1.15}
\begin{tabularx}{\textwidth}{>{\raggedright\arraybackslash}p{0.22\textwidth} >{\raggedright\arraybackslash}p{0.36\textwidth} Y}
\textbf{Lag} & \textbf{Indhold} & \textbf{Faglig brug} \\
\hline
iOS-app & SwiftUI-features, RootView, AppState, domain models, repository protocols, Supabase repositories, support services og designsystem. & Viser at UI, state og dataadgang er opdelt i lag frem for direkte databasekald fra views. \\
watchOS-app & WatchRootView, WatchWorkoutStore og WatchAuthSyncBridge. & Dokumenterer watchOS som companion til dagens træning og hurtig logging. \\
Extensions & Live Activity widget, Live Activity intents og watch widget/complication. & Underbygger lav-friktionsstatus under træning uden fuld appåbning. \\
Supabase & Auth, Postgres med RLS/views/RPC og Storage. & Viser backendens rolle som auth-, data- og sikkerhedsgrænse. \\
Edge Functions & Admin, trainer-client, program/import/analyse og billing functions. & Afgrænser privilegerede operationer til server-side runtime. \\
Apple services & StoreKit 2, WatchConnectivity, ActivityKit/WidgetKit og HealthKit. & Placering af platformservices og afgrænsning af integrationer. \\
\end{tabularx}
\caption{Komponentoverblik for Træningsmesters arkitektur.}
\label{tab:tm-container-architecture}
\end{table}

\begin{table}[htbp]
\small
\renewcommand{\arraystretch}{1.15}
\begin{tabularx}{\textwidth}{>{\raggedright\arraybackslash}p{0.24\textwidth} >{\raggedright\arraybackslash}p{0.34\textwidth} Y}
\textbf{Procespunkt} & \textbf{Kondenseret procesindhold} & \textbf{Faglig brug} \\
\hline
Start og session & Brugeren starter workout; AppState validerer auth og profile mode. & Viser at træningsflowet er brugerbundet og sessionafhængigt. \\
Tracker-valg & Gateway afgør, om tracker er aktiveret. & Gør tracker-off completion til en legitim vej gennem flowet. \\
Tracker-route & Appen åbner workoutTracker, henter øvelser og historik og synkroniserer Live Activity snapshot. & Dokumenterer sammenhæng mellem planlagt workout, faktisk logging og statusflade. \\
No-tracker-route & Appen åbner workoutTrackerWithoutLogging og går direkte mod completion. & Understøtter håndtering af afvigelser uden tvungen sætlogning. \\
Log-sæt-loop & Brugeren indtaster reps, vægt eller varighed; repository skriver trackerlog og opdaterer projection. & Viser hvordan gentagne sæt håndteres i et loop. \\
Completion & WorkoutCompletionService gemmer completion event. & Adskiller gennemført træning fra detaljeret trackerlog. \\
Progression & Aktiv cyklus kalder progression RPC efter completion. & Kobler completion til næste træningshandling og cyklusprogression. \\
Historik/Home & Historik og Home bruger completion- og trackerlogdata efter afslutning. & Viser feedback tilbage til brugerens næste app-state. \\
\end{tabularx}
\caption{Tracker- og completionproces for Træningsmester.}
\label{tab:tm-tracker-flow-bpmn}
\end{table}

\begin{table}[htbp]
\small
\renewcommand{\arraystretch}{1.15}
\begin{tabularx}{\textwidth}{>{\raggedright\arraybackslash}p{0.22\textwidth} >{\raggedright\arraybackslash}p{0.40\textwidth} Y}
\textbf{Autorisationsspor} & \textbf{Kondenseret regel} & \textbf{Faglig brug} \\
\hline
Miljø og klient & Appen loader miljø, afviser placeholder- eller service-role keys og opretter derefter anon Supabase client. & Underbygger beslutningen om, at service-role ikke må ligge i klienten. \\
Session og brugerdata & AuthenticateUser og ResolveSession er forudsætning for ReadOwnData, WriteOwnData og delt planadgang. & Viser auth som betingelse for brugerbundne læse-/skrivehandlinger. \\
Admin & CallAdminFunction kræver ResolveSession; derefter skal ValidateAdminServerSide ske, ellers afvises operationen. & Dokumenterer at admin-UI ikke er autoritativt i sig selv. \\
Trainer/import & Service-role er kun legal inside Edge Function runtime for trainer-, admin- eller importfunktioner. & Afgrænser privilegeret adgang til backend. \\
Tracker og Live Activity & WriteTrackerLog kræver session; Live Activity write kræver session og deduplicering via client action. & Understøtter idempotency og RLS i træningslogging. \\
Afvisning og audit & Cross-user access afvises, og privilegerede mutationer bør persistere audit event. & Kobler sikkerhed, governance og sporbarhed. \\
\end{tabularx}
\caption{Autorisationsregler for Træningsmester.}
\label{tab:tm-authorization-dcr}
\end{table}

\clearpage

\subsection{SQL, RLS og sikkerhed}

\begin{table}[htbp]
\small
\renewcommand{\arraystretch}{1.15}
\begin{tabularx}{\textwidth}{>{\raggedright\arraybackslash}p{0.22\textwidth} >{\raggedright\arraybackslash}p{0.36\textwidth} Y}
\textbf{Domæne} & \textbf{Centrale tabeller} & \textbf{Formål} \\
\hline
Øvelseskatalog & exercises, workout exercise user weights & Officielle/private øvelser, søgning, medier og brugerspecifik vægt. \\
Programmer & plan, workout, plan workouts, workout exercises & Programstruktur, træningsdage og øvelsesrækker. \\
Deling/social & plan share codes, shared users, invites, friend codes og friendships & Programdeling, invites og buddy-relationer. \\
Træner/klient & trainer clients og trainer client plan links & Klientrelationer og planadgang. \\
Match & match og liked/parrede views & Like-/swipe-state for øvelser. \\
Tracker/historik & trackerlog, workout completion events og active workout session state & Sætlogs, completion uden trackerlog og aktiv session. \\
Settings/admin & user settings, activity events, feedback, version history og error reports & Profiltilstand, subscription, admin, feedback og diagnostics. \\
Cykler/import & training cycles, cycle runtime-tabeller og training import jobs & Makro/meso/deload, runtime progression og importjobs. \\
\end{tabularx}
\caption{Centrale SQL-domæner i Træningsmesters datamodel.}
\label{tab:tm-sql-domain-overblik}
\end{table}

\begin{table}[htbp]
\small
\renewcommand{\arraystretch}{1.15}
\begin{tabularx}{\textwidth}{>{\raggedright\arraybackslash}p{0.24\textwidth} >{\raggedright\arraybackslash}p{0.24\textwidth} Y}
\textbf{Felt} & \textbf{Placering} & \textbf{Hvorfor feltet er relevant} \\
\hline
profile mode & User settings, completion, session og cycle flows & Skelner personlig bruger fra træner-/klientkontekst. \\
client action id & Trackerlog og workout completion events & Giver idempotency for Live Activity-, watch- og tracker-writes. \\
reverted at & Workout completion events & Gør det muligt at fortryde completion uden at slette historik. \\
superset group og position & Workout exercises & Understøtter superset-struktur i tracker og watch. \\
rest override & Workout exercises & Skelner standardpause fra bevidst pauseoverride. \\
subscription tier & User settings & Binder StoreKit-køb til appens featureadgang. \\
is admin & User settings & Bruges sammen med servervalidering, ikke som eneste adminbevis. \\
\end{tabularx}
\caption{Særligt relevante SQL-felter i Træningsmesters datamodel.}
\label{tab:tm-sql-fields}
\end{table}

\begin{table}[htbp]
\footnotesize
\renewcommand{\arraystretch}{1.12}
\begin{tabularx}{\textwidth}{>{\raggedright\arraybackslash}p{0.19\textwidth} >{\raggedright\arraybackslash}p{0.25\textwidth} >{\raggedright\arraybackslash}p{0.25\textwidth} Y}
\textbf{Område} & \textbf{Primær sikkerhedsgrænse} & \textbf{Klientbeskyttelse} & \textbf{Server/SQL-beskyttelse} \\
\hline
Auth/session & Supabase Auth og anon key & AppEnvironment afviser placeholders/service-role. & Auth JWT og auth.uid(). \\
Egne programmer & RLS på plan/workout-tabeller & Repositories kræver session og typed payloads. & Owner policies og helper functions. \\
Delte programmer & RLS og share views & UI viser kun accessible resources. & Read/edit helpers for delt planadgang. \\
Trackerlog & RLS og immutable field trigger & Tracker repositories sender brugerbundne writes. & Immutable trigger og unik client action. \\
Completion events & RLS og idempotency & Completion-service centraliserer skrivninger. & Unik user/action og reverted-at-felt. \\
Active session & RLS og forced RLS & App state og Live Activity snapshot. & Unique user/profile session og constraints. \\
Watch sync & App Group/local cache og anon session & Watch token payload model. & RLS for backend reads/writes. \\
Live Activity intents & App Group snapshot og anon Supabase client & SharedStore og projectionmodel. & RLS og idempotente writes. \\
Admin & Edge Function og SQL guard & UI viser kun adminflow efter server-backed state. & Admin guards og service-role kun i function runtime. \\
Trainer/client & Edge Function og RLS & Trainer-client-service. & Related-user policies og trainer-client function. \\
AI/import & Edge Functions & Klienten kalder ikke OpenAI direkte. & Function runtime styrer service-role/OpenAI. \\
\end{tabularx}
\caption{RLS- og sikkerhedsmatrix for Træningsmester.}
\label{tab:tm-rls-matrix}
\end{table}

\clearpage

\subsection{Beta, brugerindsigt og verifikation}

\begin{table}[htbp]
\footnotesize
\renewcommand{\arraystretch}{1.08}
\begin{tabularx}{\textwidth}{>{\raggedright\arraybackslash}p{0.18\textwidth} >{\raggedright\arraybackslash}p{0.16\textwidth} >{\raggedright\arraybackslash}p{0.18\textwidth} Y}
\textbf{Dato} & \textbf{Version/build} & \textbf{Platform} & \textbf{Faglig brug} \\
\hline
2026-03-25 22:20 UTC & Invitation & iOS & Start på iOS-beta og formel rekruttering. \\
2026-03-26 16:29 UTC & 1.0 (5) & iOS + watchOS & Første observerede build efter invitation. \\
2026-03-27 13:02 UTC & 1.0 (6) & iOS + watchOS & Iteration samme uge som betaåbning. \\
2026-03-27 13:17 UTC & 1.0 (8) & iOS + watchOS & Hurtig efterfølgende build. \\
2026-03-29 23:29 UTC & 2.0 (1) & iOS + watchOS & Skift til 2.0-sporet. \\
2026-03-30 23:07 UTC & 2.0 (2) & iOS + watchOS & Build omkring link-/opdateringsfriktion. \\
2026-03-31 00:09 UTC & 2.0 (3) & iOS + watchOS & Natlig iteration op til større releasebesked. \\
2026-03-31 10:58 UTC & 2.0 (4) & iOS + watchOS & Samme dag som katalog-, søgnings- og releasefeedback. \\
2026-03-31 19:24 UTC & 2.0 (5) & iOS + watchOS & Iterativ build efter feedback samme dag. \\
2026-04-01 18:19 UTC & 2.0 (6) & iOS + watchOS & Fortsat testdistribution efter martsfeedback. \\
2026-04-15 00:40 UTC & 2.0 (7) & iOS + watchOS & Ny beta efter april-iteration. \\
2026-04-15 01:02 UTC & 2.0 (8) & iOS + watchOS & Hurtig efterfølgende iteration. \\
2026-04-20 00:57 UTC & 2.0 (9) & iOS + watchOS & Pre-audit/releasekandidat. \\
2026-04-23 15:19 UTC & 2.0 (10) & iOS + watchOS & Sen beta med fortsat polish. \\
2026-04-23 23:37 UTC & 2.0 (11) & iOS + watchOS & Samme døgn som build 10; kort releasecyklus. \\
2026-04-25 21:01 UTC & 2.0 (12) & iOS + watchOS & Seneste observerede beta-build i den gennemgåede betaperiode. \\
\end{tabularx}
\caption{TestFlight-tidslinje for Træningsmester. Tabellen dokumenterer distribution og iteration, ikke sundhedseffekt.}
\label{tab:tm-testflight-timeline}
\end{table}

\begin{table}[htbp]
\footnotesize
\renewcommand{\arraystretch}{1.12}
\begin{tabularx}{\textwidth}{>{\raggedright\arraybackslash}p{0.14\textwidth} >{\raggedright\arraybackslash}p{0.18\textwidth} >{\raggedright\arraybackslash}p{0.32\textwidth} Y}
\textbf{ID} & \textbf{Tema} & \textbf{Anonymiseret observation} & \textbf{Designimplikation} \\
\hline
B01 & Friktion & Lock screen/Live Activity blev bekræftet som fungerende i praktisk brug. & Prioriter hurtig status under træning uden fuld appåbning. \\
B02 & Mediekvalitet & Dialogen pegede på ønske om rigtige øvelsesvideoer frem for generiske klip. & Øvelsesmedier skal være troværdige og kropsligt genkendelige. \\
B05 & Betaproces & Build 4-12 blev distribueret mellem 2026-03-31 og 2026-04-25. & Dokumenterer iterationstempo, men ikke brugerudbytte alene. \\
B10 & Rekruttering & iOS-kravet udelukkede mindst én interesseret deltager. & Betaen bør beskrives som iOS-biased convenience sample. \\
B11 & Installationsfriktion & TestFlight-link, Apple review-status og manuel opdatering skabte forvirring. & Beta-onboarding skal vise build/version og fallback-instruktion tydeligt. \\
B12 & Katalog & Risiko for dubletter og lav datakvalitet ved fri øvelsesoprettelse. & Private øvelser og offentligt katalog bør adskilles med review-flow. \\
B13 & Søgning & Tester kunne ikke finde almindelige øvelser eller varianter i betaen. & Søgning skal understøtte synonymer, varianter og øvelsesfamilier. \\
B15 & Enheder & Tester ønskede kg/lbs, fordi nogle maskiner bruger lbs. & Tilbyd global enhed, per-øvelse override og logkonvertering. \\
B16 & Watch & Apple Watch blev positioneret som telefonfri træningsflade med stabilitetsforbehold. & Prioriter oversigt, hurtig logging og robust session-sync. \\
B18 & Noter & Feedback pegede på behov for noter til teknik, udstyr eller sædeindstilling. & Øvelsesnoter bør behandles som træningskontekst, ikke kun fritekst. \\
\end{tabularx}
\caption{Udvalgte anonymiserede betaobservationer. Direkte citater og personhenførbare detaljer er fravalgt.}
\label{tab:tm-beta-feedback}
\end{table}

\begin{table}[htbp]
\small
\renewcommand{\arraystretch}{1.15}
\begin{tabularx}{\textwidth}{>{\raggedright\arraybackslash}p{0.26\textwidth} >{\raggedright\arraybackslash}p{0.25\textwidth} Y}
\textbf{Teknisk påstand} & \textbf{Dokumenteret status} & \textbf{Begrænsning} \\
\hline
Appen kan bygges som iOS-klient & Bestået & Ikke dokumenteret som uafhængig reproducerbar build. \\
Central forretningslogik er unit-testet & Bestået & Dækker unit-lag, ikke alle live backendscenarier. \\
Supabase anon read path fungerer & Delvist bestået & Tests med admin- og ikke-admin-legitimationsoplysninger blev sprunget over. \\
watchOS targetet kan bygges & Bestået & Runtime-forløb er ikke dokumenteret visuelt i bilaget. \\
Live Activity targetet kan bygges & Bestået & Dynamic Island-forløb er ikke dokumenteret visuelt i bilaget. \\
Aktuel iOS UI er dokumenteret visuelt & Bestået & Bygger på de aktuelle screenshots, der er indsat i bilaget. \\
Data-/sikkerhedsmodellen er dokumenteret & Stærk dokumentation & Fuld security proof kræver flere integrationstests. \\
\end{tabularx}
\caption{Verifikationsspor for tekniske claims.}
\label{tab:tm-test-traceability}
\end{table}

\clearpage
\fussy
